%% RetailPulse -- IEEE-style project report (IEEEtran, conference).
%%
%% Build with the project helper so the output is reproducible:
%%     python -m docs.generate_ieee_report
%% which sets SOURCE_DATE_EPOCH and runs latexmk. See docs/ARCHITECTURE.md.
\documentclass[conference]{IEEEtran}

% Reproducible output: pdfTeX stamps a timestamp and a trailer id into every
% PDF unless told not to. Guarded so the file still builds on engines that do
% not define these primitives.
\ifdefined\pdfinfoomitdate \pdfinfoomitdate=1 \fi
\ifdefined\pdftrailerid \pdftrailerid{} \fi
\ifdefined\pdfsuppressptexinfo \pdfsuppressptexinfo=-1 \fi

\usepackage{cite}
\usepackage{graphicx}
\usepackage{url}
\usepackage[T1]{fontenc}

\begin{document}

\title{RetailPulse: A Compact End-to-End Analytics Pipeline as a Vehicle for
Data-Quality Engineering}

\author{\IEEEauthorblockN{Saurabh Mastud}
\IEEEauthorblockA{Independent project\\
\url{https://github.com/SaurabhMastud/retailpulse}}}

\maketitle

\begin{abstract}
Tutorials on the modern data stack typically demonstrate one link of the chain
in isolation, which leaves the failure modes that only appear between
components unexercised. This paper reports on RetailPulse, a deliberately small
but complete batch analytics pipeline --- synthetic event generation, a raw
landing zone, validating ingestion into an embedded warehouse, SQL
transformation as code, workflow orchestration, and a read-only dashboard ---
built over seven working sessions and 55 commits as a single-developer
engineering case study. The system itself is unremarkable by design; the
contribution of this report is an account of its data-quality strategy and of
the defects that strategy caught. We classify the project's 107 automated
checks into five layers spanning unit tests, warehouse constraint tests,
cross-source agreement tests, and pipeline-health guardrails, and we document
six concrete defects that reached a committed state and were subsequently
detected by a specific layer, together with their root causes. A recurring
pattern emerges: every defect that survived review was one where a check
compared an artifact against itself rather than against an independent
reference, or where no check read the produced artifact back at all. We
conclude with the limitations of an $n{=}1$ case study and the threats to
validity of the evidence presented.
\end{abstract}

\begin{IEEEkeywords}
data engineering, data quality, ELT, dbt, DuckDB, workflow orchestration,
software testing, reproducible builds
\end{IEEEkeywords}

\section{Introduction}

The components of a contemporary analytics stack are individually well
documented. An event source feeds a landing zone; an ingestion process loads
raw records into a warehouse; transformation-as-code tools such as dbt
\cite{dbt} build layered models; an orchestrator such as Apache Airflow
\cite{airflow} schedules the whole thing; and a thin consumption layer renders
the result. Material explaining any one of these links is plentiful.

What is scarcer is an account of the defects that arise \emph{between} the
links. A generator that is correct in isolation can still produce data that
makes a downstream aggregate ambiguous. A warehouse test that passes on every
batch can be structurally incapable of failing. A guard that is green in a
developer's working copy can be guaranteed to fail on a fresh checkout. These
are integration-level faults, and a single-component demonstration cannot
surface them.

RetailPulse was built to put the entire chain under test at a scale small
enough that the whole thing stays comprehensible. The domain is synthetic
retail commerce --- page views, cart additions and purchases --- chosen because
the resulting analytics (revenue per day, funnel conversion, best-selling
products) are intuitive enough that a wrong answer is recognisable as wrong.

This report makes three claims, each supported by artifacts in the public
repository rather than by measurement on a production system:

\begin{enumerate}
\item A complete, honestly tested pipeline of this shape fits in roughly 2{,}000
  lines of Python and 11 SQL files, which is small enough for one developer to
  hold in mind and therefore small enough to reason about.
\item Automated checks are usefully classified by \emph{what independent
  reference they compare against}, not by the tool that runs them. We give a
  five-layer taxonomy on that basis.
\item Of the six defects documented here, every one was admitted by a check
  that lacked such an independent reference --- either comparing an artifact to
  itself, or never reading the produced artifact back.
\end{enumerate}

Section~\ref{sec:related} situates the work. Section~\ref{sec:arch} describes
the architecture and Section~\ref{sec:impl} the implementation.
Section~\ref{sec:dq} presents the data-quality taxonomy and
Section~\ref{sec:eval} the defects it caught. Section~\ref{sec:discussion}
discusses lessons and threats to validity.

\section{Background and Related Work}
\label{sec:related}

Dimensional modelling, and in particular the separation of conformed dimensions
from fact tables, is long-established practice \cite{kimball}. RetailPulse
follows it in miniature: a product dimension is maintained independently of the
event stream, and the marts join the dimension rather than trusting attributes
carried on each event. Section~\ref{sec:eval} shows that this separation is
what made one class of defect detectable at all.

Automated data-quality verification at scale is addressed by systems such as
Deequ, which expresses declarative constraints over datasets and verifies them
as part of a pipeline \cite{deequ}. The present work operates several orders of
magnitude below that setting and uses dbt's built-in schema tests plus
hand-written singular tests instead, but adopts the same central premise: data
constraints belong in version control and run on every load.

The warehouse is DuckDB, an embeddable analytical database \cite{duckdb}, which
removes server administration from a single-developer project while keeping the
SQL dialect and the dbt integration realistic. The consumption layer uses
Streamlit \cite{streamlit}. Synthetic user activity is drawn from a
Zipf-weighted popularity distribution \cite{zipf}, so that a minority of
simulated users generate a disproportionate share of events, as observed in
real traffic. Report reproducibility follows the \texttt{SOURCE\_DATE\_EPOCH}
convention of the Reproducible Builds project \cite{sde}.

\section{System Architecture}
\label{sec:arch}

The pipeline is a batch ELT chain with an immutable landing zone. Stages are
listed in Table~\ref{tab:components}.

\begin{table}[t]
\caption{Pipeline stages and their implementations}
\label{tab:components}
\centering
\begin{tabular}{|l|l|}
\hline
\textbf{Stage} & \textbf{Implementation} \\
\hline
Event generation & Python, JSON Lines output \\
Landing zone & Local filesystem, append-only \\
Ingestion & Python validation + DuckDB upsert \\
Warehouse & DuckDB, single file \\
Transformation & dbt: 6 models, 1 seed, 3 sources \\
Orchestration & Airflow DAG, \texttt{@daily} \\
Consumption & Streamlit, read-only \\
\hline
\end{tabular}
\end{table}

Data moves in one direction. The generator writes one immutable JSON Lines file
per batch into \texttt{data/raw/}, named by a zero-padded UTC timestamp so that
lexicographic order is chronological order. Ingestion validates each record,
routes rejects to a quarantine table with the reason attached, and upserts
accepted records on \texttt{event\_id}. A \texttt{pipeline\_runs} audit table
records per-batch counts of records read, loaded, duplicated and rejected. dbt
then builds a staging layer of typed views and three marts at daily grain:
\texttt{daily\_revenue}, \texttt{funnel\_conversion} and
\texttt{top\_products}.

Two architectural decisions are load-bearing for the rest of this report.

\textbf{Idempotent ingestion.} Because ingestion upserts on a natural key, the
same landing file can be loaded any number of times without changing the
result. This was introduced for retry safety, but it is also what made
landing-zone replay --- rebuilding a deleted warehouse from raw files --- cost
no new loading logic at all.

\textbf{Logic outside the orchestrator.} The pipeline stages are plain Python
functions in one module; the Airflow DAG only wires them together. The
immediate motivation was that Airflow has no native Windows support and the
project was developed on Windows, so a DAG holding the logic would have been
neither runnable nor testable. The durable benefit is that every stage is
callable from a test without an orchestrator present.

\section{Implementation}
\label{sec:impl}

\subsection{Event generation}
Events are emitted in browsing \emph{sessions} rather than as independent
draws: one to four events share a user, a session identifier and clustered
timestamps. This is a requirement of the funnel mart, which counts
view-to-cart-to-purchase transitions per visit; independent draws would make
the funnel's denominators meaningless. Session start times are spread over a
configurable window (14 days by default) so that date-grained marts have more
than one date to group by. Simulated users have one or two preferred
categories, and activity levels follow a Zipf weighting \cite{zipf}.

\subsection{Ingestion and quarantine}
Validation checks required fields, known event types, and price/quantity
sanity. Invalid records are not discarded: they are written to a quarantine
table with the raw payload and the rejection reason, so that a batch which
silently rejects a large fraction of its rows leaves evidence behind. The
distinction between a genuine insert and a skipped duplicate is obtained from
the \texttt{RETURNING} clause of the upsert rather than inferred from row
counts, which also makes within-batch duplicates detectable.

\subsection{Transformation}
The staging layer types and renames raw columns and nothing else. The marts sit
at the finest useful grain --- per date, and per (date, product) for
\texttt{top\_products} --- on the principle that an all-time aggregate can
answer exactly one question, whereas any window can be rolled up downstream
from a daily grain. A session is attributed to the date of its first event, so
a visit crossing midnight is counted once.

The product dimension is a dbt seed \emph{exported from} the generator's
catalogue rather than derived from the events. This distinction matters: a
dimension built as \texttt{SELECT DISTINCT product\_id, product\_category FROM
stg\_events} agrees with the events by construction and can therefore never
contradict them.

\subsection{Consumption and reporting}
The dashboard separates layout from warehouse reads, so the logic that can
actually be wrong --- window arithmetic, rate derivation, roll-ups --- is
tested against a built warehouse while the rendering gets a smoke test through
Streamlit's own test runner. Project reports are generated from the
architecture document by a Python renderer and committed; the PDF creation
timestamp is pinned so that a diff on the committed artifact indicates a real
content change.

\section{Data-Quality Strategy}
\label{sec:dq}

The project carries 107 automated checks: 74 Python tests (one skipped where
Airflow is absent) and 33 dbt data tests. Counting them is less informative
than classifying them by the reference they compare against, which is what
determines the faults they are capable of detecting.
Table~\ref{tab:layers} gives the taxonomy.

\begin{table}[t]
\caption{Check layers by independent reference}
\label{tab:layers}
\centering
\begin{tabular}{|l|l|}
\hline
\textbf{Layer} & \textbf{Compared against} \\
\hline
Unit & An expected value stated in the test \\
Schema constraint & A declared invariant (unique, not null) \\
Cross-source & A second, independently produced artifact \\
Pipeline health & The pipeline's own audit and clock \\
Artifact read-back & The produced output, re-parsed \\
\hline
\end{tabular}
\end{table}

The third and fifth layers are the ones most often missing in practice, and
they are where this project's real defects were found. A cross-source check
requires that two artifacts be produced by genuinely different paths --- hence
exporting the product dimension from the generator's catalogue, which lets a
referential test fail when an event names a product the catalogue has never
heard of. An artifact read-back check requires parsing the thing that was
actually produced, rather than asserting that producing it did not raise.

Pipeline-health checks are deliberately split in two. A volume check asserts
that the \emph{latest} audit row did not read a batch and then load, duplicate
and reject nothing, which is the signature of a stalled loader as distinct from
a legitimately empty source. A freshness check asserts that the newest event
timestamp has not fallen further behind the present than the generator's own
window, which catches the opposite failure: a warehouse that is voluminous but
no longer advancing.

\section{Evaluation: Defects Detected}
\label{sec:eval}

This is an engineering case study, not a controlled experiment. We make no
performance claims; the warehouse holds a few thousand rows and rebuilds in
under a second, so throughput figures would measure nothing of interest.
Instead we report, as the primary evidence, the defects that reached a
committed state and were later detected, with the layer that caught each one.
Table~\ref{tab:defects} summarises them.

\begin{table}[t]
\caption{Defects detected, by the layer that caught them}
\label{tab:defects}
\centering
\begin{tabular}{|l|l|}
\hline
\textbf{Defect} & \textbf{Caught by} \\
\hline
Per-batch product catalogue & Schema constraint \\
Events-only category check & Cross-source (replacement) \\
Seed drift on fresh checkout & Cross-source, in a clone \\
Missing seed in test fixture & Cross-source (cascade) \\
Stale warehouse & Pipeline health \\
Mangled report identifiers & Artifact read-back \\
\hline
\end{tabular}
\end{table}

\textbf{Per-batch product catalogue.} The generator's catalogue was seeded with
the per-run random seed, so each batch invented its own product-to-category
mapping. Once the product mart was re-grained to (date, product), a single
product split into duplicate rows. The composite-key constraint failed on the
first multi-batch run. Root cause: reference data generated with the same seed
as the random data that references it.

\textbf{Events-only category check.} A test asserted that each product mapped
to exactly one category, comparing events against other events. It could not
fail in the worst case: if every event for a product agreed on the \emph{same
wrong} category, the test passed. It was replaced by a comparison against the
exported dimension, which is strictly stronger. This defect was in the test
suite, not the pipeline --- a check that cannot fail is a defect in its own
right.

\textbf{Seed drift on fresh checkout.} A guard asserted that the committed seed
file byte-matched a fresh export. It passed in the working copy and would have
failed on every fresh clone, because the version control system was configured
to rewrite line endings on checkout while the exporter writes Unix endings. It
was found by cloning the repository, not by running the suite.

\textbf{Missing seed in test fixture.} A dashboard test fixture built models
without first loading the seed. The omission was harmless until a mart began
referencing the seed, at which point it failed. The ordering rule had been
duplicated across four call sites and was already wrong in one of them within
hours of being introduced; it was consolidated into a single function.

\textbf{Stale warehouse.} The freshness check failed against a local warehouse
that had not been loaded for approximately six weeks of elapsed time, and
failed again two weeks later. On both occasions the resolution was to run the
pipeline, not to widen the threshold. A check that can be satisfied by relaxing
its own bound is not a check.

\textbf{Mangled report identifiers.} The report renderer stripped underscores
as inline emphasis markup. In a document composed largely of
\texttt{snake\_case} identifiers, this corrupted every one of them across six
pages, for several commits. The pre-existing test asserted only that a non-empty
PDF had been written, which such a report satisfies. Extracting the rendered
text and asserting on it surfaced the fault immediately. Root cause: markup
rules written for prose applied to a document that is mostly identifiers.

Four of the six defects were admitted by a check that compared an artifact to
itself or to nothing; the remaining two were admitted because no check read the
produced artifact back. None were found by inspection.

\section{Discussion}
\label{sec:discussion}

\subsection{Lessons}
Three generalisations survive the specifics. First, the question to ask of any
data test is what independent artifact it compares against; if the answer is
``the same table'', it verifies internal consistency and not correctness.
Second, a generated artifact that is also committed must be deterministic, or
its version history becomes noise --- this applies equally to an exported CSV
and to a rendered PDF. Third, a passing local suite is not evidence that a
project runs: two of the six defects above were invisible until the repository
was cloned and executed from a bare checkout.

\subsection{Threats to validity}
Several limitations bound these conclusions.

\emph{Construct validity.} The defect count is the number \emph{detected}, not
the number present. Defects that no layer was capable of catching are by
definition absent from Table~\ref{tab:defects}, and their number is unknown.

\emph{Internal validity.} The author designed the system, wrote the checks and
classified the defects. The taxonomy was constructed after the defects were
found, so its fit to them is not independent evidence of its generality.

\emph{External validity.} This is a single project at small scale with
synthetic data, built by one developer. Concerns that dominate production
systems --- late-arriving data, schema evolution, multi-writer concurrency,
cost --- are absent by construction, and nothing here should be read as
evidence about them.

\emph{Unexercised components.} The Airflow DAG has never been executed. The
development platform has no native Airflow support, so the DAG is validated
structurally: its task identifiers, dependency order, and the continued
existence of every function it calls are asserted at the abstract-syntax-tree
level. Everything it orchestrates does run, through the same functions, via the
command-line entry point. The dashboard likewise runs only locally. There is no
continuous integration, so all checks are enforced by local discipline rather
than by the repository.

\section{Conclusion and Future Work}
\label{sec:conclusion}

RetailPulse demonstrates that a complete batch analytics chain, honestly
tested, is tractable at a scale one developer can keep in mind: approximately
2{,}000 lines of Python, 11 SQL files, 6 dbt models and a seed, built across 7
sessions and 55 commits. Its more transferable output is the finding that the
useful axis for classifying data-quality checks is the independence of their
reference, and the observation that every defect which survived review in this
project was admitted by a check lacking one.

The most valuable next step is continuous integration: it is the single
documented limitation with a concrete fix, and it would have caught the
checkout-dependent defect on the first clean build. A secondary item is to
resolve the dashboard's warehouse path at call time, so that the render checks
execute in environments where no warehouse has yet been built rather than
skipping.

\begin{thebibliography}{9}

\bibitem{kimball}
R.~Kimball and M.~Ross, \emph{The Data Warehouse Toolkit: The Definitive Guide
to Dimensional Modeling}, 3rd~ed. Indianapolis, IN, USA: Wiley, 2013.

\bibitem{duckdb}
M.~Raasveldt and H.~M\"{u}hleisen, ``DuckDB: an embeddable analytical
database,'' in \emph{Proc. ACM SIGMOD Int. Conf. Management of Data}, 2019,
pp.~1981--1984.

\bibitem{deequ}
S.~Schelter, D.~Lange, P.~Schmidt, M.~Celikel, F.~Biessmann, and A.~Grafberger,
``Automating large-scale data quality verification,'' \emph{Proc. VLDB
Endowment}, vol.~11, no.~12, pp.~1781--1794, 2018.

\bibitem{dbt}
dbt Labs, ``dbt documentation.'' [Online]. Available:
\url{https://docs.getdbt.com}

\bibitem{airflow}
Apache Software Foundation, ``Apache Airflow documentation.'' [Online].
Available: \url{https://airflow.apache.org/docs/}

\bibitem{streamlit}
Streamlit Inc., ``Streamlit documentation.'' [Online]. Available:
\url{https://docs.streamlit.io}

\bibitem{zipf}
G.~K.~Zipf, \emph{Human Behavior and the Principle of Least Effort}. Cambridge,
MA, USA: Addison-Wesley, 1949.

\bibitem{sde}
Reproducible Builds Project, ``SOURCE\_DATE\_EPOCH specification.'' [Online].
Available: \url{https://reproducible-builds.org/docs/source-date-epoch/}

\end{thebibliography}

\end{document}
